% ============================================================
% Chapter 4 -- The Gaussian chain by belief propagation
%
% Introduces only the factor-graph and sum-product facts the thesis needs,
% then the chain sweeps, the Gaussian closure (= Kalman/RTS) and the measured
% price of locality.
% ============================================================

\chapter{The Gaussian Chain by Belief Propagation}
\label{ch:gaussian-bp}\label{ch:graphs}\label{ch:graphical}

Chapter~\ref{ch:gaussian} obtained the joint score by inverting the dense
$\nsites\times\nsites$ matrix $\Sigt$, at $O(\nsites^3)$. That cost belongs
to the route, not to the problem: the posterior precision $J$ of
Proposition~\ref{prop:g-posterior} is tridiagonal, so a structured solve
delivers the same answer in $O(\nsites)$. This chapter obtains the score by
the message-passing form of that solve, sum--product on the factor graph of
the posterior, generic and exact on any chain; for the Gaussian chain a
message collapses to two numbers and the recursion is the Kalman filter
followed by the Rauch--Tung--Striebel smoother. The chapter closes with the
cost of locality: what truncating the evidence range does to the score.

\section{Factor graphs and the sum--product algorithm}
\label{sec:factor-graphs}\label{sec:sum-product}\label{sec:mrf}

A factor graph makes a factorisation graphical
\citep{kschischang2001factor}: given $p(x) = Z^{-1}\prod_a \psi_a(x_{\partial a})$
it is the bipartite graph with a variable node per $x_i$, a factor node per
$\psi_a$, and an edge whenever $x_i \in x_{\partial a}$. A connected factor
graph with one fewer edge than nodes is a \emph{tree}. Sum--product belief
propagation (BP) passes messages along every edge
\citep{pearl1988probabilistic, mezard2009information}: variable-to-factor
messages multiply the incoming factor messages from all \emph{other}
neighbours, and factor-to-variable messages integrate the factor against the
incoming variable messages,
\begin{equation}
  \nu_{i \to a}(x_i) \propto
    \prod_{b \in \partial i \setminus a} \hat\nu_{b \to i}(x_i),
  \qquad
  \hat\nu_{a \to i}(x_i) \propto
    \int \psi_a(x_{\partial a})
    \prod_{j \in \partial a \setminus i} \nu_{j \to a}(x_j)\;
    \dd x_{\partial a \setminus i},
  \label{eq:bp-rules}
\end{equation}
with the belief at node $i$ proportional to $\prod_{a \in \partial i}
\hat\nu_{a \to i}(x_i)$ (Figure~\ref{fig:bp-messages}). The ``exclude the
recipient'' rule is the heart of the algorithm: the message $i \to a$
summarises everything node $i$ knows except what it learned from $a$, so no
evidence is counted twice along a path.

\begin{figure}[t]\centering
\includegraphics[width=0.78\linewidth]{fig_mezard_bp_messages.png}
\caption[The two sum--product update rules]{The two halves of
\eqref{eq:bp-rules}. \emph{Left:} the message from variable $j$ to factor
$a$, a function of the incoming factor messages from every other factor
neighbour $b \ne a$. \emph{Right:} the message from factor $a$ to variable
$j$, a function of the incoming variable messages from every other variable
neighbour $k \ne j$. The original carries iteration superscripts; on a tree
one inward and one outward sweep suffice. Reproduced from
\citet{mezard2009information}, their Figure~14.3.}
\label{fig:bp-messages}
\end{figure}

\begin{theorem}[Tree exactness of sum--product]\label{thm:tree-exact}
On a tree-shaped factor graph, one inward and one outward pass of
\eqref{eq:bp-rules} compute the exact node and factor marginals of $p$
\citep[Theorem~14.1]{mezard2009information}. The beliefs are unique; the
unnormalised messages are defined up to positive constants.
\end{theorem}

\begin{proof}[Proof sketch]
Each edge of a tree splits the graph into two disjoint halves, so the message
crossing it is exactly the partial sum over the half behind it, and no
evidence reaches a node twice. Induction from the leaves inward terminates
with every marginal having absorbed each factor once.
\end{proof}

On graphs with cycles the same updates, iterated, define loopy BP, which is
approximate \citep{weiss2001correctness}; none of that is needed here. Bayes'
rule multiplies the chain prior \eqref{eq:chain-factorisation} by the
per-frame likelihoods of \eqref{eq:ou-channel}, and nothing in the likelihood
couples two frames, so for the Gaussian chain
\begin{equation}
  p(a \,|\, x) \;\propto\;
  \underbrace{\psi_{\mathrm{pr}}(a_0)}_{\Nd(a_0;\,0,\,1)}
  \prod_{k=0}^{\nsites-2}
  \underbrace{\psi^{(k)}_{\mathrm{tr}}(a_k, a_{k+1})}
            _{\Nd(a_{k+1};\,\alpha a_k,\,\ivar)}
  \prod_{k=0}^{\nsites-1}
  \underbrace{\psi^{(k)}_{\mathrm{ob}}(a_k)}
            _{\Nd(x_k;\,e^{-t}a_k,\,\Deltat)} ,
  \label{eq:bp-posterior}
\end{equation}
a hidden Markov model \citep{rabiner1989tutorial} with the data $x_k$ absorbed
into unary observation factors. Its factor graph has $\nsites$ variable
nodes, $2\nsites$ factor nodes and $3\nsites - 1$ edges: a tree, a spine of
variables linked by transition factors with one evidence leaf per node
(Figure~\ref{fig:posterior-factor-graph}). Conditioning on the data did not
change the graph of the prior, because independent per-node observation
re-weights a node potential and never joins two spine nodes; correlated
observation noise would add edges and break the construction. Nothing in this
paragraph used Gaussianity: the posterior of \emph{any} chain prior under
coordinatewise noise is a tree, and Theorem~\ref{thm:tree-exact} applies to
it.

\begin{figure}[t]\centering
\begin{tikzpicture}[
  latent/.style={circle, draw, minimum size=7mm, inner sep=1pt},
  obs/.style={circle, draw, fill=black!10, minimum size=7mm, inner sep=1pt},
  fac/.style={rectangle, draw, fill=black, minimum size=2.2mm, inner sep=0pt},
  node distance=9mm, scale=0.92, every node/.style={transform shape}]
\foreach \i in {1,2,3,4}{
  \node[latent] (a\i) at (2.1*\i, 0) {$a_{\i}$};
  \node[fac] (f\i) at (2.1*\i, -1.0) {};
  \node[obs] (x\i) at (2.1*\i, -2.0) {$x_{\i}$};
  \draw (a\i) -- (f\i) -- (x\i);
}
\foreach \i/\j in {1/2,2/3,3/4}{
  \node[fac] (w\i) at (2.1*\i+1.05, 0) {};
  \draw (a\i) -- (w\i) -- (a\j);
}
\node[above=1.5mm of w1] {\small $\kernel(a_{k}\mid a_{k-1})$};
\node[right=2mm of f4] {\small $\like(x_k\mid a_k)$};
\draw[-{Stealth}, thick] (1.15,1.35) -- (2.55,1.35) node[midway,above] {\small $\lmsg$};
\draw[-{Stealth}, thick] (9.35,1.35) -- (7.95,1.35) node[midway,above] {\small $\rmsg$};
\end{tikzpicture}
\caption[Posterior factor graph of a noised Markov chain]{The posterior factor
graph of a Markov chain observed through the coordinatewise channel. Open
circles are latent variables, shaded circles observations, squares factors.
Each transition factor joins two adjacent latent variables; each observation
enters through a unary factor attached to a single latent variable. The
posterior graph is therefore still a chain, Theorem~\ref{thm:tree-exact}
applies, and the arrows show the two message sweeps.}
\label{fig:posterior-factor-graph}
\end{figure}

Tree exactness is a property of the graph and says nothing about whether the
messages can be \emph{represented}: for discrete variables a message is a
finite vector, for continuous variables a density, and the recursion is made
computable by closure in a finite-dimensional family
(Section~\ref{sec:bp-gaussian}), by a numerical grid (the reference solver
of Chapter~\ref{ch:laplace}), or by projection onto a tractable family at a
cost Chapter~\ref{ch:laplace} measures.

\section{The chain sweeps and their cost}
\label{sec:bp-chain}\label{sec:bp-sweep}

On the tree of Figure~\ref{fig:posterior-factor-graph} the schedule of
\eqref{eq:bp-rules} collapses to one forward and one backward pass.

\begin{proposition}[Chain sweeps]\label{prop:bp-convA}
Define
\begin{align}
  \mu_{\to 0}(a_0) &:= \psi_{\mathrm{pr}}(a_0),
  &
  \mu_{\to k+1}(a_{k+1}) &:= \int
    \psi^{(k)}_{\mathrm{tr}}(a_k, a_{k+1})\,
    \psi^{(k)}_{\mathrm{ob}}(a_k)\,
    \mu_{\to k}(a_k)\; \dd a_k ,
  \label{eq:bp-fwd}\\
  \mu_{\leftarrow \nsites-1}(a_{\nsites-1}) &:= 1,
  &
  \mu_{\leftarrow k-1}(a_{k-1}) &:= \int
    \psi^{(k-1)}_{\mathrm{tr}}(a_{k-1}, a_k)\,
    \psi^{(k)}_{\mathrm{ob}}(a_k)\,
    \mu_{\leftarrow k}(a_k)\; \dd a_k .
  \label{eq:bp-bwd}
\end{align}
Then for every $k$,
\begin{equation}
  p(a_k \,|\, x) \;\propto\;
  \mu_{\to k}(a_k)\;
  \psi^{(k)}_{\mathrm{ob}}(a_k)\;
  \mu_{\leftarrow k}(a_k) ,
  \label{eq:bp-combine}
\end{equation}
where $\mu_{\to k}$ depends on $x_0,\dots,x_{k-1}$ only, $\mu_{\leftarrow k}$
on $x_{k+1},\dots,x_{\nsites-1}$ only, and the local evidence enters once.
\end{proposition}

\begin{proof}
Fix $k$ and split the factors of \eqref{eq:bp-posterior} into those left of
$a_k$, the local observation, and those right of $a_k$. With $a_k$ held
fixed the two sides share no variable, so integrating out everything else
factorises into a left integral, which is $\mu_{\to k}(a_k)$ by definition,
and a right integral, which is $\mu_{\leftarrow k}(a_k)$. Peeling off the last
integration variable of the left integral gives \eqref{eq:bp-fwd}; the first
of the right integral gives \eqref{eq:bp-bwd}.
\end{proof}

For a chain of $K$-state variables every message is a vector in $\R^K$, one
transition update is a $K \times K$ matrix--vector product, and a
forward--backward sweep sends $O(\nsites)$ of them: time $O(\nsites K^2)$
and space $O(\nsites K)$, linear in the sequence length.

\section{The Gaussian closure: Kalman filtering and RTS smoothing}
\label{sec:bp-gaussian}

Two facts close the recursion algebraically: a product of Gaussians in the
same variable is Gaussian, and an affine image of a Gaussian plus independent
Gaussian noise is Gaussian. Write an unnormalised scalar Gaussian in natural
form, $\mu(a) \propto \exp(-\tfrac12\lambda a^2 + r a)$, with precision
$\lambda$ and field $r$, so mean $r/\lambda$ and variance $1/\lambda$. The
observation factor is itself Gaussian in $a_k$,
$\psi^{(k)}_{\mathrm{ob}}(a_k) \propto \exp\big(-\tfrac12\lambda_{\mathrm{ob}} a_k^2
+ r_{\mathrm{ob},k} a_k\big)$ with
\begin{equation}
  \lambda_{\mathrm{ob}} = \frac{e^{-2t}}{\Deltat},
  \qquad
  r_{\mathrm{ob},k} = \frac{e^{-t}x_k}{\Deltat}:
  \label{eq:bp-observation-parameters}
\end{equation}
a pseudo-observation at the deconvolved location $e^{t}x_k$ with precision
$e^{-2t}/\Deltat$.

\begin{theorem}[Gaussian closure of BP on the chain]\label{thm:bp-closure}
For the Gaussian AR(1) chain observed through the OU channel every
sum--product message and every belief is Gaussian. With $(\lambda, r)$ the
natural parameters of the message arriving at site $k$ and
$(\bar\lambda, \bar r) = (\lambda + \lambda_{\mathrm{ob}},\, r + r_{\mathrm{ob},k})$
after absorbing the local observation, the sweeps are
\begin{equation}
  \lambda_{\to k+1} = \frac{\bar\lambda}{\alpha^2 + \ivar\bar\lambda},
  \quad
  r_{\to k+1} = \frac{\alpha\bar r}{\alpha^2 + \ivar\bar\lambda};
  \qquad
  \lambda_{\leftarrow k-1} = \frac{\alpha^2\bar\lambda}{1 + \ivar\bar\lambda},
  \quad
  r_{\leftarrow k-1} = \frac{\alpha\bar r}{1 + \ivar\bar\lambda},
  \label{eq:bp-forward-precision}
\end{equation}
and the smoothing belief at site $k$ adds the three natural parameters,
\begin{equation}
  \lambda_k^{\mathrm{post}} = \lambda_{\to k} + \lambda_{\mathrm{ob}} + \lambda_{\leftarrow k},
  \qquad
  r_k^{\mathrm{post}} = r_{\to k} + r_{\mathrm{ob},k} + r_{\leftarrow k},
  \qquad
  \E[a_k \mid x] = \frac{r_k^{\mathrm{post}}}{\lambda_k^{\mathrm{post}}},
  \label{eq:bp-belief-mean}
\end{equation}
from which \eqref{eq:g-tweedie} gives the score. Each update is $O(1)$, so all
posterior marginals, and hence the complete joint score, are obtained exactly
in $O(\nsites)$ time and space.
\end{theorem}

\begin{proof}
Multiplying Gaussian factors adds natural parameters, which gives
$(\bar\lambda, \bar r)$. Forward, $a_k \sim \Nd(\bar r/\bar\lambda, 1/\bar\lambda)$
pushed through $a_{k+1} = \alpha a_k + \eta_k$ is
$\Nd(\alpha\bar r/\bar\lambda,\ \ivar + \alpha^2/\bar\lambda)$, whose natural
parameters are the first pair in \eqref{eq:bp-forward-precision}. Backward,
$\int \Nd(a_k; \alpha a_{k-1}, \ivar)\exp(-\tfrac12\bar\lambda a_k^2 + \bar r a_k)\,\dd a_k$
is a Gaussian function of $a_{k-1}$ with the second pair. The belief
\eqref{eq:bp-combine} is a product of three Gaussians. No approximation is
introduced at any step: the family is closed under the two operations the
algorithm uses, and the graph is a tree.
\end{proof}

Run in mean--variance arithmetic the forward sweep is the Kalman filter,
predict $(m, v) \mapsto (\alpha m,\ \alpha^2 v + \ivar)$ followed by a
precision-weighted update with the observation, and combining it with the
backward sweep is the Rauch--Tung--Striebel smoother, update for update
\citep{kalman1960new, rauch1965maximum}: the two classical literatures
compute one object. The closure is not the Gaussian approximation of TAP or
approximate message passing \citep{thouless1977solution, donoho2009message},
which compress many weak incoming messages into Gaussian updates with an
Onsager correction in a large-system limit; here the messages are Gaussian
because multiplication and marginalisation are exact operations on the
family, on a tree of finite degree. Tree exactness holds for any innovation
law while closure holds only because the model is Gaussian, and
Chapter~\ref{ch:laplace} is about the second property failing while the first
survives. Numerically, the audit suite runs BP against the matrix route
\eqref{eq:g-score-matrix} over $80$ configurations
($\nsites \in \{2,3,5,8,16\}$, $\alpha \in \{0.2, 0.5, 0.9, -0.4\}$,
$t \in \{0.05, 0.3, 1, 3\}$): maximum disagreement $1.3 \times 10^{-15}$ on
posterior means and $1.2 \times 10^{-14}$ on scores, double-precision
rounding accumulated along the recursion, reproduced by an independent
implementation written from the recursion alone (Appendix~\ref{app:repro}).
One quantitative detail from an $\nsites = 3$ instance: at
$x = (1.2, -0.4, 0.8)$ with $\alpha = 0.9$, $t = 0.05$ the score component
$S_1$ is $+5.71$, fourteen times the $-x_1 = +0.40$ a per-frame marginal
score would give: the neighbours, not the frame itself, set the size of the
correction.

\section{The price of locality}
\label{sec:bp-locality}\label{sec:g-bulk}\label{sec:g-locality}\label{sec:bp-truncation}

Belief propagation on the chain is exact but not local: the message reaching
frame $k$ has travelled the whole sequence, and $\Qt$ is dense for every
$t > 0$. How much of the score survives if evidence is allowed to travel only
a few frames? There are two ways to go local.

\emph{Truncate the answer.} Keep the entries of the exact operator within a
band and use what is left,
\begin{equation}
  \widehat S^{(b)}(x,t) \;=\; -\,\Qt^{(b)} x,
  \qquad
  \big(\Qt^{(b)}\big)_{ij} = (\Qt)_{ij}\,\mathbf{1}\{|i-j| \le b\};
  \label{eq:bp-banded}
\end{equation}
$b = 1$ is the tridiagonal, Markov, truncation of the score.
\emph{Truncate the algorithm.} Let the messages of
Proposition~\ref{prop:bp-convA} travel at most $r$ hops before frame $k$ is
read off, replacing the message beyond the cut by the stationary prior
marginal. Because a contiguous window of a stationary AR(1) chain is itself a
stationary AR(1) chain, this is exact inference on the window,
\begin{equation}
  m^{(r)}_k \;=\; \E\big[a_k \,\big|\, x_{k-r}, \dots, x_{k+r}\big],
  \qquad
  \widehat S^{(r)}_k(x,t) \;=\; \frac{e^{-t} m^{(r)}_k - x_k}{\Deltat}.
  \label{eq:bp-windowed}
\end{equation}
Both estimators are linear, $\widehat S = -Mx$, so their error against
$S = -\Qt x$ can be averaged over the data law $x \sim \Nd(0,\Sigt)$ without
sampling. Using $\E[x^\T A x] = \operatorname{tr}(A\Sigt)$ and
$\operatorname{tr}(\Qt\Sigt\Qt) = \operatorname{tr}\Qt$, the relative RMS
score error is the closed form
\begin{equation}
  \boxed{\;
  \varepsilon(M)
  :=
  \sqrt{\frac{\E\|(M-\Qt)x\|^2}{\E\|\Qt x\|^2}}
  =
  \sqrt{
    \frac{\operatorname{tr}\!\big[(M-\Qt)\,\Sigt\,(M-\Qt)^{\T}\big]}
         {\operatorname{tr}\Qt}
  } ,
  \;}
  \label{eq:bp-relerr}
\end{equation}
a population error, not a sampled estimate: $\varepsilon = 0.1$ means the RMS
approximation error is a tenth of the RMS size of the exact score.
Table~\ref{tab:g-truncation} and Figure~\ref{fig:truncation} evaluate it for
both families.

\begin{center}
\begin{minipage}{\linewidth}\centering
\captionof{table}[The price of locality]{Relative RMS score error
\eqref{eq:bp-relerr} of the two truncations ($\nsites = 40$, $\alpha = 0.8$);
every entry is closed-form. The last two rows are derived: the fraction of the
no-neighbour error removed by admitting a single hop, and the penalty for
truncating the answer rather than the algorithm at equal budget.}
\label{tab:g-truncation}
\small
\begin{tabular}{@{}lcccc@{}}
\toprule
relative RMS score error & $t = 0.05$ & $t = 0.3$ & $t = 1$ & $t = 3$ \\
\midrule
no neighbours, $r = 0$        & $0.815$ & $0.567$ & $0.194$ & $0.0045$ \\
one hop, $r = 1$              & $0.123$ & $0.235$ & $0.130$ & $0.0035$ \\
two hops, $r = 2$             & $0.020$ & $0.107$ & $0.089$ & $0.0028$ \\
\addlinespace
banded answer, $b = 1$        & $0.210$ & $0.301$ & $0.135$ & $0.0035$ \\
banded answer, $b = 2$        & $0.032$ & $0.134$ & $0.092$ & $0.0028$ \\
\midrule
removed by the first hop      & $85\%$  & $58\%$  & $33\%$  & $21\%$ \\
penalty, $b{=}1$ over $r{=}1$ & $1.71\times$ & $1.28\times$ & $1.03\times$ & $1.00\times$ \\
\bottomrule
\end{tabular}
\end{minipage}
\end{center}

The first row is the cost of treating the frames as independent: at
$t = 0.05$ the isolated-frame score is wrong by $81.5\%$ in relative RMS, so
almost all of the score at small diffusion times \emph{is} inter-frame
structure, and the error does not vanish as $t \to 0$ because the $r = 0$
estimator converges to $-x$ while the exact score converges to $-\Qz x$. A
single nearest-neighbour exchange, the Markov band and nothing more, removes
$85\%$ of it, and two hops remove $98\%$. The removed fraction falls with
diffusion time, but the first row shows why the band is not mattering less:
by $t = 3$ the entire inter-frame content of the score is worth $0.0045$ in
relative error, so nothing is left for any hop to recover. The band carries
the largest share of the coupling exactly where the coupling is large. At
equal budget, truncating the inference beats truncating the score matrix at
small $t$ ($0.123$ against $0.210$ at $t = 0.05$), and the gap closes as $t$
grows. The error of the radius-$r$ estimator decays geometrically in $r$ at
a rate that worsens with $t$; the closed-form influence-decay rate behind
that geometry is derived in Chapter~\ref{ch:em-results}, where it predicts
the best receptive field of a learned score head.

\begin{figure}[t]\centering
\includegraphics[width=0.92\linewidth]{fig_truncation.pdf}
\caption[The price of locality]{The two truncations ($\nsites = 40$,
$\alpha = 0.8$; every value closed-form by \eqref{eq:bp-relerr}).
\emph{Left:} relative RMS score error against diffusion time, for message
range $r$ (solid) and matrix band $b$ (dashed). For $r \ge 1$ both are worst
at intermediate $t$ and vanish at both ends, at $t \to 0$ because the clean
score is tridiagonal; the no-neighbour estimator is the exception, since the
clean score of frame $k$ genuinely uses its neighbours whenever
$\alpha \ne 0$. \emph{Right:} the same error against the locality budget at
three diffusion times: geometric decay at a rate that worsens with $t$, with
the banded answer (squares) above the windowed inference (circles) at every
budget.}
\label{fig:truncation}
\end{figure}
